\section{Trade-off Analysis (Prototype Comparison)}

This section compares the prototypes on operational complexity, performance, reliability, and developer velocity.
It does \emph{not} pick a winner yet; it defines how we will decide.

\subsection{Evaluation Criteria}
\begin{itemize}
  \item \textbf{Simplicity:} how many moving parts must be deployed and understood.
  \item \textbf{Coupling:} how strongly services depend on each other's uptime and latency.
  \item \textbf{Backpressure handling:} behavior during bursts and downstream slowdowns.
  \item \textbf{Latency:} time-to-ticket and time-to-enrichment.
  \item \textbf{Reliability:} replayability, safe retries, and failure isolation.
  \item \textbf{Operability:} debugging, observability, and day-2 operations.
  \item \textbf{Cost:} infrastructure + operational overhead.
  \item \textbf{Data retention:} raw event storage, reprocessing, and auditing.
\end{itemize}

\subsection{Comparison Summary}

\begin{longtable}{@{}p{0.18\linewidth}p{0.26\linewidth}p{0.26\linewidth}p{0.26\linewidth}@{}}
\toprule
\textbf{Dimension} & \textbf{Prototype A (Direct REST)} & \textbf{Prototype B (Broker)} & \textbf{Prototype C (Log Store + Broker)} \\
\midrule
Infra components & Lowest (3 services + DB) & Medium (adds broker + DLQ) & Highest (adds broker + log store) \\
Coupling & High (sync dependencies) & Medium (async ingestion; optional sync NLP) & Low/Medium (async by default) \\
Burst handling & Weak unless watcher buffers & Strong (broker absorbs spikes) & Strong (broker + raw storage separation) \\
Time-to-ticket & Low (direct call path) & Medium/Low (consume lag under load) & Medium/Low (consume lag; pointer events) \\
Time-to-enrichment & Low if NLP stable & Medium (async option) & Medium (async default; fetch from store) \\
Replay/reprocess & Hard (needs custom logs) & Good (replay from broker) & Best (replay pointers + reread raw) \\
Payload size limits & Tight (HTTP/request size) & Medium (broker limits) & Best (broker carries pointers only) \\
Debugging & Simple trace path & Needs queue tooling & Needs queue + store tooling \\
Cost & Lowest & Medium & Highest \\
\bottomrule
\end{longtable}

\subsection{Common Failure Modes and Mitigations}
\textbf{Applies to all prototypes:}
\begin{itemize}
  \item \textbf{Duplicate events:} use \texttt{event\_id} for idempotency and deduplication.
  \item \textbf{Hot sources:} add per-source rate limiting and sampling policies.
  \item \textbf{Noisy events:} implement suppression windows and grouping keys (e.g., \texttt{source + category + signature}).
\end{itemize}

\textbf{Prototype A specific:}
\begin{itemize}
  \item \textbf{Downstream slowness cascades:} watcher and core must enforce timeouts and fail fast.
  \item \textbf{NLP outage blocks requests:} create tickets without enrichment and enrich later via retry queue.
\end{itemize}

\textbf{Prototype B specific:}
\begin{itemize}
  \item \textbf{Poison messages:} route to DLQ after bounded retries; build replay tooling.
  \item \textbf{Exactly-once assumptions:} assume at-least-once; make consumers idempotent.
\end{itemize}

\textbf{Prototype C specific:}
\begin{itemize}
  \item \textbf{Pointer refers to missing raw payload:} enforce write-before-publish and include \texttt{content\_hash}.
  \item \textbf{Store growth:} define retention policies and index lifecycle management early.
\end{itemize}

\subsection{Decision Rubric (How We Will Choose)}
We pick the prototype that best fits the first target environment and workload.
Use these questions as the selection checklist:
\begin{enumerate}
  \item What is the peak ingestion rate and expected burst size?
  \item Is ``ticket created within X seconds'' a hard SLO?
  \item Is enrichment required for correctness, or can it be eventually consistent?
  \item Do we need replay/reprocessing for audits or model improvements?
  \item How much infrastructure is acceptable for the MVP (Kafka, log store, Supabase)?
  \item What is the operational maturity of the team running it?
\end{enumerate}
